\documentclass[11pt]{article}

\usepackage[margin=1in]{geometry}
\usepackage{microtype}
\usepackage{parskip}
\usepackage[hidelinks]{hyperref}
\usepackage{titlesec}
\usepackage{enumitem}
\usepackage{booktabs}
\usepackage{amssymb}

\setlist{nosep, leftmargin=1.5em, topsep=2pt}
\titleformat{\section}{\large\bfseries}{\thesection.}{0.6em}{}
\titleformat{\subsection}{\normalsize\bfseries}{\thesubsection}{0.6em}{}
\titlespacing*{\section}{0pt}{1.6em}{0.5em}
\titlespacing*{\subsection}{0pt}{1.1em}{0.3em}

\hypersetup{
  pdftitle={An automation shock in a small economy},
  pdfauthor={Robin Langer},
  pdfborder={0 0 0}
}

\newcommand{\hyp}[1]{\medskip\noindent\textit{Hypothesis.}\ #1\medskip}

\title{\textbf{An automation shock in a small economy}}
\author{Robin Langer}
\date{5 November 2026}

\begin{document}
\maketitle

\begin{abstract}
\noindent
I propose to build a small, closed, agent-based economy in which a single
real phenomenon --- the \emph{automation of production} --- is the only
internal shock, and to ask whether that one change is sufficient to produce
the two movements Irving Fisher separated: a fall in the price of goods, and
a debt-deflation in the price of assets. The economy is stock-flow consistent:
every unit of money is somebody's asset and somebody else's liability, created
only when a bank lends and destroyed only when a loan is repaid. Every agent
is boundedly rational and sees only its own books; no agent forecasts the
macro-economy. The claims of interest --- deflation, structural unemployment,
a debt-deflation spiral, and a grim mortality-driven floor beneath it --- are
things the model should be \emph{able to produce from micro-rules}, never
assumptions written into its equations. This document specifies the entities,
the markets, the money, and the clock in enough detail to implement; every
design question is resolved, and what remains is numeric calibration.
\end{abstract}

{\small\tableofcontents}

%====================================================================
\section{The proposal}
%====================================================================

\subsection{The question}

The nightmare scenario is simple to state and hard to produce honestly. A
generation of machines arrives that makes \emph{more} goods with \emph{fewer}
workers. Supply rises and employment falls at the same instant; the newly
unemployed stop spending, so demand falls while supply climbs. Prices fall ---
deflation. But debt is a fixed sum of money, and the interest on it must still
be paid out of a money supply that is now shrinking. Repaying yesterday's loan
in today's dearer money raises the \emph{real} burden of every debt, which
pushes households and firms to sell assets to raise cash, which deflates asset
prices further: a debt-deflation spiral. High real rates kill the incentive to
start or grow a business, investment collapses, and the economy settles into a
depressed state.

The model's job is not to assert this story but to see whether it
\emph{emerges}. I will give agents local rules --- buy food before anything
else, lower your price when goods go unsold, refuse a loan that looks unsafe
--- and watch whether the macro-pattern above appears without anyone being told
to produce it. The single internal cause is one declining parameter: the number
of operators a machine of a given type and generation requires.

\subsection{Modelling stance}

Three commitments shape every later choice.

\textit{Agent-based.} The unit of the model is the individual household or
firm, acting on its own accounts, the prices posted so far, the calendar, and
its own recent sales. There is no representative agent and no central
auctioneer.

\textit{Stock-flow consistent (SFC).} Money is not a free-floating price level.
It is created when a bank writes a new deposit against collateral, and
destroyed when that loan is repaid; a loan between two individuals only
\emph{moves} money that already exists. Every flow leaves an accounting trace
in two places. A slump shrinks the money supply because banks make fewer new
loans, not because anyone's cash is annihilated.

\textit{Boundedly rational; the market prices risk, not the agent.} No agent
holds a model of the economy or forecasts the cycle. A bank is a mechanical
secured lender that checks a collateral ratio. A private lender applies a local
heuristic to a borrower's disclosed accounts (does current revenue cover the
monthly payment with margin?). Risk is priced by the \emph{market} --- a
borrower posts terms, lenders accept or decline, and a firm that cannot place
its debt re-posts at a higher rate next turn --- never by a forecast inside
anyone's head.

%====================================================================
\section{The world}
%====================================================================

\subsection{Entities and resources}

The economy contains four kinds of agent: \textbf{individuals} (households, some
of whom are also entrepreneurs or mine owners), one or two \textbf{banks}, and
a \textbf{government}. Everything they trade falls into a small set of
resource classes.

\begin{itemize}
  \item \textbf{Consumer goods and services}, bought to raise welfare: staple
  food (survival), luxury food and private healthcare (general health),
  nightclubs and bigger housing (happiness), and private tuition (education).
  \item \textbf{Real estate} --- houses of various sizes --- which is both a
  place to live (a requirement, below) and an illiquid asset.
  \item \textbf{Liquid financial assets}: private debt, government debt, and
  bank shares. These trade every month.
  \item \textbf{Managed productive assets}: \emph{mines} (mining land plus its
  finite stock of raw material) and \emph{businesses} (a bundle of machines and
  the employees who operate them).
  \item \textbf{Machines}, which come in distinct \emph{types} and, within a
  type, successive \emph{generations} (levels). The generation is the
  automation dial of \S\ref{sec:bijection}.
  \item \textbf{Raw materials} --- a common and an advanced type --- both
  consumed to build machines, the advanced one increasingly at higher
  generations (\S\ref{sec:raw}).
\end{itemize}

A household's welfare is three stocks --- health, happiness, education --- and
its survival turns on a fourth necessity, food, bought before anything else.
Individuals pay no tax; firms do.

\subsection{The agent archetypes}

The same individual can wear several hats; the archetypes below are roles, not
castes, and the model lets agents move between them.

\begin{itemize}
  \item \textbf{An ordinary individual} holds a job (or does not), lives in a
  house it owns or rents, carries welfare stocks, and each month must feed
  itself before it spends, invests, or saves.
  \item \textbf{An entrepreneur} is an individual who also owns a business
  outright --- machines plus employees --- and must meet the firm's
  obligations (wages, tax, debt service) before taking any profit home.
  \item \textbf{A tech entrepreneur} runs the special business that builds
  \emph{machines} out of raw materials, employing engineers. This is the only
  agent that can introduce a new machine generation, so it is the supply valve
  of automation itself.
  \item \textbf{A mine owner} holds mining land, must post a selling price and
  quantity for its raw material each month, may sell the mine on the secondary
  market, and may deposit it into a bank in exchange for bank shares.
  \item \textbf{The indigenous population} are the initial mine owners: at the
  start they own all mining land and the mines, which makes them the economy's
  first rentier-financier class (\S\ref{sec:seed}).
  \item \textbf{The government} taxes sales, employs public doctors, teachers
  and bureaucrats, runs the public hospital and university, borrows when short,
  and provides health and education in kind --- but never cash, housing, or
  food.
  \item \textbf{The bank} is a mechanical secured lender whose capital base is
  deposited mining land, who may lend up to a multiple of that base, and who
  pays its depositor-shareholders a yearly dividend.
\end{itemize}

%====================================================================
\section{How a person lives}
%====================================================================

\subsection{Welfare: health, happiness, education, and the hard zeros}
\label{sec:welfare}

Each lever of consumption does exactly one thing; there are no cross-effects to
blur the accounting. The luxury supermarket and the private hospital raise
\emph{general health}; the nightclub and a bigger house raise \emph{happiness};
teachers (public or private) raise \emph{education}; the ordinary supermarket
buys staple food, which is survival rather than welfare. Doctors are a special
case: they do not raise general health at all --- they \emph{cure diseases}
(\S\ref{sec:health-market}).

Two couplings turn these tidy stocks into a trap. First, \textbf{happiness caps
health}: there is a ceiling on how healthy an unhappy agent can be, so even a
single-minded health-maximiser must buy some happiness to lift the ceiling
before more medicine or luxury food does anything. Second, there are two
\textbf{hard-zero} states that pin health at zero and hold it there ---
\emph{homelessness} and an \emph{uncured disease}. While pinned, no amount of
luxury food helps; the condition must be resolved first (get housed; or buy the
specialist cure). And because employers prefer healthier applicants
(\S\ref{sec:job-market}), an agent at health zero is also unhireable. Job loss
then lowers income, which costs the agent its housing or its health, which makes
it harder to be hired --- a micro-level poverty spiral that is the household
mirror of the macro debt-deflation spiral.

Homelessness itself does not kill; it sends health to zero but the agent lives,
and public services still reach it, so the trap is escapable. What kills is a
sustained inability to buy food, under a \textbf{three-turn rule}. An agent that
cannot afford food does not die at once: it enters starvation, and the engine
automatically force-sells its \emph{liquid} assets --- cash first, then bank
shares, debt holdings, and directly-held mine land --- at a below-market distress
price to feed it. Only after three consecutive foodless turns, once nothing
liquid remains to burn, does it die. Because houses can be sold only once a year
(\S\ref{sec:time}) and are not force-sold here, an asset-rich but cash-poor owner
whose only wealth is an unsellable house still dies this way --- death by
\emph{illiquidity}, not insolvency (\S\ref{sec:realestate-sale}). The same
three-turn rule governs a newborn that cannot be fed (\S\ref{sec:birthdeath}).

\subsection{Careers, education levels, and the machine-operator bijection}
\label{sec:bijection}

Every productive job is the operation of a machine, and machine types stand in
\textbf{bijection} with worker classes: a bureaucrat runs the admin machine, a
doctor the medical machine, a teacher the teaching machine, and retail and
hospitality staff theirs, while an \emph{engineer} runs the meta-machine that
builds machines. The thing that automates is itself a machine with operators who
can, in turn, be automated. Mining is the exception: extraction is
\emph{labour-free}, a mine producing raw material directly from its land and
finite stock with no operators at all (\S\ref{sec:raw}).

Alongside these productive workers, every business also employs
\textbf{bureaucrats} who run its admin machine and \emph{support} the productive
staff in a fixed ratio --- $\beta$ bureaucrats per productive worker, about one
to one at generation~0 (a hospital with twenty doctors runs twenty bureaucrats).
Because the admin machine automates like any other, a higher
generation needs fewer bureaucrats per unit of throughput, so $\beta$ falls: the
back office sheds workers exactly as the shop floor does.

Education is \textbf{per career}, and within a career it is a \emph{level}. Your
level in a career decides which \emph{generation} of that career's machine you
may operate: level~0 runs generation~0, higher levels run higher generations.
An ``advanced engineer'' is therefore not a separate profession but simply a
high-level engineer --- the one who can build the next machine generation.

The friction that makes automation bite is that \textbf{switching careers resets
education to zero}. Your level belongs to one career and does not transfer; a
bureaucrat displaced by automation who retrains as an engineer starts again at
level~0. This is what turns automation unemployment \emph{structural} rather
than merely cyclical: the displaced cannot glide sideways into the jobs that are
growing. Which career an agent educates toward is its own decision, made by
looking at last year's job offerings and training toward whatever was most in
demand --- adaptive, backward-looking expectations rather than a forecast.

\subsection{An ordinary individual's turn}
\label{sec:individual-turn}

Each month, an individual who is neither an entrepreneur nor a mine owner acts
in a strict, necessity-first order:

\begin{enumerate}
  \item \textbf{Get paid.} Collect all income due: wages (if employed), the
  monthly repayment on any debt it holds, mine-land yield, and --- once a year,
  as an upfront lump --- rent from any property let out. The government also
  offers, this turn, a rationed chance at a free public level-up or disease cure
  (\S\ref{sec:education-market}, \S\ref{sec:health-market}).
  \item \textbf{Buy food, or starve.} The first act of every turn. An agent that
  cannot buy food enters starvation; the engine force-sells its liquid assets
  below market to feed it, and death follows only after three consecutive
  foodless turns (\S\ref{sec:welfare}).
  \item \textbf{Buy welfare and education} with what remains --- health,
  happiness, education, in any amount supply allows, on top of whatever the
  government provided free.
  \item \textbf{Invest} with what is still left, in the monthly capital market:
  mine land and debt (\S\ref{sec:capital-market}). Houses are \emph{not} bought
  here --- they trade only on the annual real-estate day.
  \item \textbf{Save} whatever is left as cash.
\end{enumerate}

This necessity-first ordering is one of three parallel turn-orders --- the
household's here, the entrepreneur's in \S\ref{sec:entrepreneur-turn}, the
government's in \S\ref{sec:government} --- all sharing one principle: in a
downturn each agent is squeezed from the bottom up, and the order decides what
gets cut first.

%====================================================================
\section{The markets}
%====================================================================

\subsection{One protocol, many venues}
\label{sec:protocol}

Every market in the economy runs the \emph{same} disequilibrium protocol, so no
market ever clears instantaneously and gluts and shortages are both possible and
informative. Each turn:

\begin{enumerate}
  \item \textbf{Sellers post} their prices and quantities first.
  \item \textbf{Buyers decide} what and how much to buy.
  \item If demand exceeds supply, the good \textbf{sells out} and a
  \textbf{random lottery} rations who is served --- rationing by chance, not by
  willingness to pay.
  \item Sellers \textbf{reprice from last turn's feedback}: raise what sold out,
  cut what went unsold, and undercut or match what competitors charged.
\end{enumerate}

One detail matters for the crash: the protocol must be \textbf{symmetric}. Just
as oversubscribed demand raises prices, an oversupply of sellers with too few
buyers must force asking prices \emph{down}. This falling-price branch is the
asset-collapse channel; without it the debt-deflation spiral cannot run.

\subsection{The consumer market (monthly)}

Food, entertainment, private health, and private education all trade monthly
under the common protocol: providers advertise, buyers choose, goods can sell
out and trigger a lottery, and prices adapt. Housing is the one consumer good
that does \emph{not} trade here --- it has its own annual day
(\S\ref{sec:realestate-sale}, \S\ref{sec:rental}).

\subsection{The education market: public and private}
\label{sec:education-market}

Education raises a citizen's level in one chosen career. Teachers are generic
--- a teacher can teach any subject --- so what is produced is set by which
subjects the government and universities choose to staff.

\emph{Public} education is a policy lever. Each turn the government picks which
teacher-types to employ, setting the mix of education produced. Citizens apply
for the units needed to reach their next level; the government fills this demand
in a priority order by housing status --- \textbf{homeless first, then renters,
then owners} --- with a \textbf{lottery} inside each tier, and allocation is
\textbf{all-or-nothing}: a winner completes a full level-up, a loser gets
nothing. So free education is real but rationed, and only in the types the
government chose to produce. \emph{Private} tuition buys the specific career-type
you want outright, escaping both the mix and the lottery.

Because a career switch resets education (\S\ref{sec:bijection}), the
government's \emph{mix} choice shapes the future labour force: keep training
bureaucrats when the economy needs engineers and you mis-train the population
for obsolete work. Steering the mix toward growing careers is the main public
tool for absorbing the automation shock.

\subsection{The health market: disease, doctors, public and private}
\label{sec:health-market}

General health is bought at the luxury supermarket; \emph{disease} is a
separate, catastrophic state. There are five disease types, and each turn
disease strikes citizens at random, sending the victim's health to a hard zero
until cured. A cure requires buying ten units of the \emph{matching} specialist
doctor (five diseases to five specialties), at a public or private hospital.

Doctor supply flows from the education system --- a generic teacher can train
any specialty --- and the government chooses its doctor-type mix each turn, just
as it chooses its teacher mix. Free public healthcare then works identically to
public education: apply, be served homeless-first with a lottery inside each
tier, all-or-nothing (the full ten units and a cure, or nothing). Private
healthcare buys the needed specialist outright. The poor therefore depend on the
homeless-first lottery while the rich self-insure privately, and a government
whose doctor mix mismatches the diseases actually circulating leaves some
citizens curable only by paying.

\subsection{The job market (annual, on the special day)}
\label{sec:job-market}

Hiring happens once a year, in the job step of the annual special day
(\S\ref{sec:time}) --- after real-estate sales, before the rental market; between
special days there is no hiring, though firms may shed workers in any month. The incumbent in a post has
first preference and may keep it at the same salary without reapplying. An agent
may also apply elsewhere, provided it meets the job's \textbf{minimum education}
--- the hard gate on who may apply (\S\ref{sec:bijection}). Among non-incumbent
applicants the \textbf{healthiest wins}: health is both a welfare stock and a
hiring edge, which is why a hard-zero agent is unemployable. A worker raises its
wage only by winning a higher-paid listing on job day, not by continuous
bargaining. A laid-off worker cannot be rehired until the next annual job day, so
unemployment can last up to twelve months --- making job loss a durable,
dangerous event that savings must bridge. Because rent is paid a year upfront, a
worker can lose its \emph{job} mid-year but not its \emph{house}: unemployment is
a monthly risk, homelessness only an annual one.

\subsection{The real-estate sale market (annual) and the illiquidity death}
\label{sec:realestate-sale}

On the special day, real-estate \textbf{sales come early} --- before the job
market, since a sale needs no job (\S\ref{sec:time}). Houses are bought and sold
only on this day --- when multiple buyers compete the winner is chosen at random
--- so real estate is illiquid between days. That illiquidity is precisely what
enables the asset-rich, cash-poor \textbf{illiquidity death} of
\S\ref{sec:welfare}: an owner who runs out of cash in month three cannot reach
its wealth until the special day months away, and houses are not among the liquid
assets the engine force-sells to feed a starving agent. Mine land and debt, by
contrast, trade monthly, so a house is the asset that can kill its owner by being
unsellable. Once sales settle, each owner --- the buyer who just acquired it, or
the owner who kept it --- \textbf{decides whether to occupy it or rent it out};
only properties offered for rent enter the rental market, which runs later, after
the job step.

\subsection{The rental market (annual, last on the special day)}
\label{sec:rental}

Renting is settled last on the special day --- after sales, the occupy-or-rent
decision, and the job market (\S\ref{sec:time}), so that applicants' salaries are
known. The owner re-advertises the rent; the sitting tenant has first preference
to renew (unless the home was sold out from under them); among other applicants
the owner prefers the one with the \textbf{highest salary or job title}. Rent is
paid \textbf{upfront for the whole year}: a tenant who cannot cover a full year
does not secure the home, and an agent left with no dwelling has its health go to
zero (homelessness). Because the year is paid in advance, housing is then secure
until the next special day --- there is no mid-year eviction. Ranking applicants
by income means that in a downturn landlords swap unemployed tenants for
higher-paid ones, so the weakest lose housing first.

\subsection{The machine market (annual, on job day)}

Tech companies sell machines to the other businesses, competing on price:
several producers list prices, businesses place orders, and machines can sell
out and push producers to raise prices next turn. Each machine carries both a
purchase cost and a maintenance cost. Machines trade on \emph{job day}, so a firm
reconfigures its capital and its workforce in the same annual act --- a machine
purchase is inseparable from the hiring or firing of its operators, and hiring is
annual.

\subsection{The raw-materials market (monthly)}
\label{sec:raw}

Mines sell raw material to tech companies every month under the common protocol.
Extraction is \textbf{labour-free}: a mine employs no one and produces raw
material directly from its land and finite stock. Each mine advertises a price
and may sell any quantity up to its remaining stock --- as much as it likes, at
whatever price it posts; offering less than the stock is withholding, which props
up the price and slows depletion. A buyer's order may carry a \textbf{roll-over}
instruction: if it is not filled, the order persists and re-tries at a higher
price up to a reservation ceiling, so a buyer that truly needs material bids
itself up within its limit. Raw material is \textbf{finite} --- each mine holds
a stock sized to it, extraction depletes it, and the total across all mines is
calibrated to be just enough to reach full automation. As mines deplete and post
higher prices, input costs to the tech sector drift up: a slow long-run drag.

There are two raw-material types --- a \textbf{common} input and an
\textbf{advanced} one, one per mine type --- and a machine's \emph{bill of
materials} depends on its generation: every machine needs the common input in
bulk plus some of the advanced one, and the advanced requirement rises with each
generation. The advanced material is therefore a second automation bottleneck
beside the engineers: as its mine depletes, each new generation costs more of a
tightening resource --- a material brake on automation that runs alongside the
labour brake.

\subsection{The capital and secondary market (monthly)}
\label{sec:capital-market}

Three liquid assets trade monthly under the common protocol --- \textbf{mine
land}, \textbf{debt}, and \textbf{bank shares} --- and for debt, primary
issuance and secondary resale are the \emph{same} market. There is only one debt
instrument (\S\ref{sec:private-debt}); a holder reselling it lists a new price
and the buyer, who sees the underlying borrower's current disclosure and the
instrument's remaining payments, accepts or declines. No new claim is created on
a resale; an existing claim changes hands. Because debt can always be resold, it
carries a market price that tracks the borrower's health: as a borrower's
accounts deteriorate, the resale value of its debt falls and holders mark it
down and dump it \emph{before} any default --- a debt-repricing channel that
feeds the asset collapse directly, not only through defaults. Participation on
the lender side is open to any agent with enough savings; it is gated by a
savings threshold, not by class.

%====================================================================
\section{Money, finance, and the firm}
%====================================================================

\subsection{Where money comes from}
\label{sec:money}

The monetary spine is simple and strict. A \textbf{bank loan creates money}: the
bank writes a new deposit from nothing against collateral. \textbf{Repayment
destroys money}: the deposit is extinguished as the loan is paid down. A
\textbf{private loan creates no money}: it only moves an existing deposit from
lender to borrower. It follows that a boom, with banks lending freely, expands
the money supply, and a bust contracts it --- not because anyone's cash is
burned, but because new lending dries up and existing loans are repaid faster
than replaced. A default, by contrast, mainly destroys \emph{bank equity}
rather than money: the deposit the borrower already spent is still in
circulation. This distinction --- repayment shrinks money, default shrinks the
bank --- is what makes the deflation and the banking fragility two separate
stories that happen to coincide in a crash.

\subsection{Mortgages and bank loans}
\label{sec:mortgage}

To borrow from a bank you must pledge \textbf{collateral}: real estate or mining
land only, never machines.\footnote{Machines are wrong-way collateral: a firm
usually fails \emph{because} its machines are an obsolete generation --- worth
least exactly when seized --- so the bank will not lend against them.} The bank
advances a lump sum and is repaid in an equal monthly stream to a deadline, the
total exceeding the lump sum by the interest. The bank is a \textbf{mechanical}
lender: it lends up to a loan-to-value fraction of the collateral, marks the
collateral to market each turn, and issues a margin call when the collateral
falls below the loan. The borrower meets the call by \textbf{selling liquid
assets first} to pay the loan down; the bank forecloses only if those cannot
cover the shortfall. It runs no cash-flow forecast and holds no model of the economy; the
loan-to-value haircut is its entire risk model. Pledged collateral carries a
\textbf{lien} that blocks sale --- you keep the use and the income of a liened
house or mine, but you cannot sell it until the loan clears, which sharpens the
illiquidity risk for an entrepreneur who pledged a home to start a firm.

\subsection{Private debt}
\label{sec:private-debt}

The direct debt market is where forward-looking credit lives, though even here
no one forecasts the macro-economy. An issuer exposes all its accounts as a
structured disclosure; it then offers an instrument --- a lump sum now, a stream
of equal monthly payments to a deadline, the total exceeding the lump sum by the
interest --- and a lender reads the disclosure and the offer and either accepts
or declines, with no negotiation. If no lender bites, the issuer re-posts next
turn at a higher rate, so a distressed firm visibly ratchets up its offered rate
to place its debt. Direct-market loans are unsecured and transferable: the
lender bears the loss if the borrower fails, which is why this debt must carry a
risk premium over collateralised bank lending.

\begin{table}[h]
\centering
\small
\begin{tabular}{@{}lll@{}}
\toprule
 & \textbf{Bank loan} & \textbf{Private loan} \\
\midrule
Security      & secured (collateral)       & unsecured \\
Credit test   & mechanical LTV gate        & lender's heuristic on disclosure \\
Transferable  & no (held to maturity)      & yes (resold on the capital market) \\
Effect on money & creates money            & moves existing money \\
Loss falls on & the bank's equity          & the individual lender \\
\bottomrule
\end{tabular}
\caption{The two credit instruments.}
\end{table}

\subsection{Government debt}

A government bond is not a special instrument; it is exactly the private-debt
instrument above, issued by the state. The government exposes its tax revenue
and its salary obligations, offers a lump sum, rate, and monthly schedule, and
savers accept or decline. Only the disclosure differs. This is what makes the
fiscal rescue of \S\ref{sec:government} endogenous: whether the state can borrow
through a downturn is a verdict of the same market that prices every other debt.

\subsection{Becoming an entrepreneur}

Entrepreneurs are not a fixed caste; an ordinary citizen can enter. The gate is
\textbf{property}: you must own real estate, because the canonical way to raise
start-up capital is to pledge that real estate to a bank and take a loan.
Renters, owning no collateral, cannot enter. Entry is therefore debt-funded ---
new firms start leveraged --- and the stake is the owner's home: because the
start-up loan is secured on it, a failure (\S\ref{sec:fold}) seizes the
\emph{pledged} property, though any other, unmortgaged property is shielded ---
and the owner keeps a food reserve and recovers at the next real-estate day by
mortgaging or selling that home.
Once running, the firm behaves like any business: it buys machines, pays wages
and taxes, and earns by selling its products. Entry makes the number of firms
endogenous, and it is the recovery channel after a crash: those who kept their
property through the bust can pledge it and start again.

\subsection{The tech entrepreneur's decisions}

The tech entrepreneur is the supply valve of automation, facing three nested
choices each period.

\begin{itemize}
  \item \textbf{Product strategy.} \emph{Stay} (same machines, same volume),
  \emph{expand} (more of the same current-generation machines, hiring more
  regular engineers --- the only growth path before automation ignites, and the
  one that \emph{adds} jobs), or \emph{upgrade} (hire a scarce advanced engineer
  to raise one product to a higher, pricier generation --- the sole way a new
  generation enters the economy).
  \item \textbf{Process and reinvestment.} When it upgrades its \emph{own}
  production machines, which lowers the operators each needs, it either
  \emph{upgrades and fires} (same output, fewer engineers, bank the saving) or
  \emph{upgrades and redeploys} (keep the engineers, make more product). Which
  it does turns on \textbf{expected demand}: a firm scales up only if it
  believes it can sell the extra output, and in a weak, deflationary state it
  sheds labour instead --- which weakens demand further. The
  displacement-versus-expansion choice is endogenous, not assumed.
  \item \textbf{Personal portfolio.} Being also a household, it may instead take
  profit home --- into happiness, health, mine land (backward integration into
  its own input supply), rental property, or others' debt.
\end{itemize}

\subsection{The entrepreneur's turn}
\label{sec:entrepreneur-turn}

Each turn an entrepreneur acts in order: pay the per-unit sales tax on what the
firm sold; service the firm's debt; decide whether to borrow (pledging real
estate or mine land, never machines); and, on job day, decide whether to buy or
upgrade machines and hire or fire operators. Only what remains after the firm's
obligations and investment may go to personal saving, outside investment, or
welfare. The ordering is the point: firm obligations precede firm investment,
which precedes personal use, so a firm in distress is cut from the bottom up and
only a genuinely insolvent one fails to meet its wages and debt.

\subsection{Firm fold and the bankruptcy waterfall}
\label{sec:fold}

When an entrepreneur can no longer meet its obligations, the firm folds. Only its
\emph{liquid} estate is marshalled --- cash savings, debt receivables, bank
shares, and directly-held mine land, all liquidated at distress prices ---
together with any specifically \emph{mortgaged} collateral the bank has a lien
on. \textbf{Unmortgaged real estate is not seized.} The proceeds meet claims in a
strict priority waterfall:

\begin{enumerate}
  \item \textbf{Wages first.} Labour is the senior claim; if the estate cannot
  pay all workers, random workers are paid until it runs out.
  \item \textbf{Tax next.} The per-unit sales tax owed on units already sold is
  paid to the government \emph{before} the bank --- a defaulting firm still owes
  the state its tax.
  \item \textbf{The bank}, recovering from the remaining liquid estate and from
  the mortgaged collateral it holds (sold slightly below market, pushing house
  prices down) and seized receivables (sold on the debt market). It does not
  reach unmortgaged property.
  \item \textbf{Other creditors} --- unsecured direct-market lenders --- take
  whatever liquid estate remains, usually nothing; the shortfall is their loss,
  money destroyed, which is why their lending carried a premium.
\end{enumerate}

\noindent\textbf{Unmortgaged property is shielded, and a food reserve is kept}:
a small survival reserve of cash --- enough to buy food until the next
real-estate day --- is carved out before the waterfall, so a fold cannot
immediately starve the owner. At that real-estate day the ex-entrepreneur
recovers by either mortgaging its unmortgaged home or selling it, turning the
shielded illiquid asset into cash. So a folded owner with an unencumbered home
carries its reserve to the next real-estate day and recovers; it dies the
three-turn starvation death only if it has neither liquid assets nor unmortgaged
property. Seizing receivables \emph{before} default closes a
laundering loophole (borrow, re-lend, fold, keep the receivable nets to zero) and
is the reason the model needs a secondary debt market at all: the bank must be
able to liquidate what it seizes.

\subsection{The mine owner's turn}

A mine owner must, each month, post a selling price and a quantity of raw
material (at most its remaining stock) and sell to tech companies under the
common protocol. It may also sell the mine itself on the monthly secondary
market, or \textbf{deposit} it into a bank. Depositing converts the mine into
\textbf{bank shares}: the physical land becomes the bank's permanent capital and
can never be withdrawn, but the depositor receives a claim on the bank's yearly
dividend, which it can sell on the monthly share market for liquidity. Whether
to keep the mine (liquid, pledgeable, sellable for emergency food money) or
deposit it (a higher but riskier, locked-up, once-a-year return) is the mine
owner's central financial choice (\S\ref{sec:bank}).

\subsection{The bank's structure and its scripted decisions}
\label{sec:bank}

The bank's capital base is \textbf{mining land}, not cash: any agent may deposit
a mine, the bank operates the deposited mines and sells their raw material, and
it may lend up to \textbf{ten times} the value of its holdings, creating money
as it lends. Depositors are the bank's shareholders; the bank pools its income
--- raw-material sales plus loan interest, minus default losses --- and pays it
out as a yearly dividend pro rata, on the special day (\S\ref{sec:time}). Because
shares trade \emph{before} the dividend that day, a share sold then carries the
dividend to its buyer. This is the one place a dividend exists, companies paying
none, and it is a strong reason to deposit.

The bank is designed to be \textbf{run-proof from the depositor side} but not
from the asset side. No depositor can withdraw capital; for liquidity it sells
its share instead, so there is no classic run. But the bank's equity still
absorbs loan losses, and because seized collateral (real estate) is illiquid,
the bank raises cash by \emph{selling its own mine-land capital} on the monthly
market. In a bust this becomes a \textbf{fire-sale amplifier}: losses force the
bank to dump mine land, mine-land prices fall, the value of its remaining base
falls, the ten-times ceiling falls, and lending contracts further --- a
mine-land debt-deflation that runs alongside the one in houses. The crash bites
the shareholders through a collapsing dividend, a falling share price, and a
shrinking base at once; cumulative losses that exhaust equity freeze credit
entirely.

Two design hazards are worth naming because they constrain the implementation.
First, a \textbf{self-valuation loophole}: if lending is sized off the posted
price of \emph{unsold} reserves, and if the deposit incentive drove everyone to
deposit so the bank \emph{is} the raw-materials market, the bank could post a
high price to inflate its own lending base and mint money against collateral it
values itself. The rule that avoids this is to \textbf{mark the deposited mines
to the independent mine-land market}: value each at its remaining stock times a
per-unit reference price drawn from recent \emph{independent} mine-land sales of
that type (adverse selection keeps the best mines independent, so that reference
market stays liquid and outside the bank's control). The bank posts
raw-material prices, never mine-land prices, so it cannot value its own
collateral; the mark reuses the same mark-to-market the bank already applies to
per-loan collateral, embeds each mine's remaining reserves automatically, and
falls in a bust as mine-land prices fall --- feeding the amplifier above. (An
optional floor at trailing realised revenue can be switched on if a mine ever
marks high while selling nothing.) Second, the
reason \textbf{not everyone deposits}: the monthly cash of independent sales
beats a yearly dividend for a household that must eat every month; deposited
capital is locked and cannot be pledged or sold for emergency food; pro-rata
pooling dilutes the best mines; and the dividend turns safe mine income into a
risky profit-share that craters in a crash. These disincentives bite hardest on
exactly the agents who would otherwise cause a monopoly --- the best mines and
the cash-poor --- so the equilibrium is a partial, adversely-selected deposit
base, which both prevents a raw-materials monopoly and preserves a competitive
price to mark the lending base against.

\subsection{The government's scripted decisions}
\label{sec:government}

The government is scripted but not omnipotent. Its income is the per-unit sales
tax plus borrowing; its spending is the wages of public doctors, teachers, and
bureaucrats, plus interest on its debt. There is \textbf{no cash welfare} ---
but public health and public education are universal in kind and reach even the
homeless, and are therefore the poverty trap's circuit-breakers. Each turn it
acts in order: receive tax; \textbf{pay its workers or fire them} (firing is its
adjustment valve, and public-sector layoffs feed straight into unemployment);
pay debt interest; issue new debt to cover any shortfall; and distribute the
health and education it produced.

Its role in the cycle is an \textbf{endogenous} stabiliser-or-amplifier. It
would rather borrow than fire, but whether it \emph{can} is decided by the debt
market: a saver reading falling tax revenue against high salary obligations may
decline or demand a punishing rate. If savers lend, the state maintains
employment but now owes interest to the wealthy --- a concentration channel; if
no one lends tolerably, it is forced into austerity --- the amplifier outcome.
The timing is cruel, because the same downturn makes the state want to borrow,
worsens its creditworthiness, and makes savers most reluctant (in deflation,
holding cash earns a positive real return). And because it too owns machines and
can automate, the public sector sheds its own workers into a system with no
welfare to catch them. The free health and education each citizen receives is
public production divided by population, so austerity, too few public machines,
or a growing population all shrink it --- weakening the circuit-breaker exactly
when it is needed most.

%====================================================================
\section{Population and time}
%====================================================================

\subsection{Birth, death, and the dark stabiliser}
\label{sec:birthdeath}

Fertility is \textbf{propertied}: only owners may have children, renters cannot,
so over generations the population descends increasingly from owners. A newborn
is endowed at birth with a fixed \textbf{required gift} --- ten units of real
estate (its own dwelling) plus a required cash gift $G$ --- and starts as a whole
new household in the lowest job tier, housed and itself fertile once it clears the
gate below. To have a child, a parent must hold, over and above its own housing,
at least \textbf{ten spare units} of real estate \emph{and} at least \textbf{$3G$
in liquid assets}, with another ten units and $3G$ for each further child;
whether an eligible parent in this voluntary band actually reproduces is governed
by its optimization strategy. Above a ceiling --- \textbf{one hundred units of
real estate and $10G$ in liquid assets} --- reproduction is \emph{mandatory}: the
very rich are forced to have a child, a strong anti-concentration bleed on the
top. On birth the engine transfers the land and raises the cash gift from the
parent's liquid assets, force-selling them at aggressively falling prices until
the gift is funded; and a newborn that cannot be fed is subject to the same
three-turn starvation rule as any agent. Real-estate accumulation and household
formation are the same process.

Death routes an estate in one of two sharply different ways. With an heir, the
firstborn inherits everything undivided --- \textbf{primogeniture}, a
concentrating force. With \emph{no} heir, the estate splits three ways: real
estate is broken into ten-unit pieces and handed to random homeless people
(instantly housing and enfranchising them); \textbf{all cash and liquid assets
escheat to the government} as revenue --- including the receivables, which the
state then collects; and the decedent's own debts are \textbf{written off}, so
its creditors lose. A childless death is thus strongly anti-concentration, the
opposite of primogeniture, and the distributional fate of the whole model hinges
on the mix of the two.

This yields the model's most distinctive emergent claim. Because a childless
death writes off the decedent's debts and de-levers the top, a \textbf{death-wave
in a crash cancels debt en masse} --- a grim jubilee by mortality. The deeper the
slump, the more deaths, the more compounding debt simply vanishes, easing the
burden on survivors and handing houses to the bottom. A second, fiscal stabiliser
rides alongside: the liquid wealth of the childless dead refills the treasury, so
a death-wave strengthens the public safety net exactly when demand for it peaks.
The spiral of \S\ref{sec:individual-turn} is therefore not purely divergent:
mortality is a genuine negative feedback, a Malthusian damper, and its losers are
disproportionately the lenders whose loan books the deaths destroy. The model
should be allowed to reach this floor honestly, through attrition, and must not be
rescued by any scripted recovery.

\hyp The economy can find a floor beneath a debt-deflation without fiscal
rescue, through mortality that de-levers it --- a smaller, poorer, de-levered
population, not a clean recovery.

\subsection{Time and scheduling}
\label{sec:time}

The clock is twelve months to a year, one turn per month. Most activity ---
food, welfare and education, the government's distribution, mine-land and debt
trading --- happens every month. But once a year a single \textbf{special day}
gathers the slow markets into one strict sequence, each step feeding the next:
(1)~\textbf{bank shares} trade, before the dividend, so a share sold now carries
this year's dividend to its buyer; (2)~the \textbf{dividend} is paid to the
post-trade holders; (3)~\textbf{real estate is bought and sold}, after which each
owner chooses to occupy or let; (4)~the \textbf{job market} opens (hiring and
machines); and (5)~the \textbf{rental market} matches tenancies --- last, because
applicants are ranked by the salary of the job they just did or did not secure.
House \emph{sales} precede the job market (a sale needs no job), while
\emph{renting} follows it. This annual rhythm has three consequences the model
should exhibit:

\begin{itemize}
  \item An \textbf{annual liquidity cliff}: a full year's rent falls due upfront
  on one day, so a worker who lost income needs a whole year's buffer to stay
  housed, and real-estate day becomes a synchronised reckoning where the
  cash-poor are evicted en masse --- a sharp annual eviction cliff rather than a
  smooth decline.
  \item \textbf{Tenant-sorting by income}: because landlords rank applicants by
  salary, a downturn lets them swap weak tenants for strong, so the weakest lose
  housing first.
  \item \textbf{Frequency separation}: housing adjusts annually while labour,
  goods, and capital adjust monthly, so housing is the slow variable, and the
  whole cycle may phase-lock to this annual forcing, with crashes clustering at
  year boundaries.
\end{itemize}

%====================================================================
\section{Behaviour, the start, and the experiment}
%====================================================================

\subsection{Optimization strategies and the social classes}

Each agent carries one objective, drawn at random at birth and not inherited,
and the dependency chains mean none is truly single-minded. A
\emph{happiness}-maximiser buys bigger houses (withdrawing them from the rental
pool) and nightclubs; a \emph{health}-maximiser must also buy happiness, because
happiness caps health; an \emph{education}-maximiser must also keep health up,
because hiring is healthiest-wins --- a three-link chain of education needing
health needing happiness; a \emph{savings}-maximiser hoards cash; an
\emph{investment}-maximiser buys income assets --- rental property, mine land,
debt. Note the opposite housing effects: happiness-maximisers withdraw real
estate from the rental market while investment-maximisers supply it, so the
strategy mix sets rental supply.

These objectives, crossed with housing and employment status, give a ladder of
\textbf{social classes} whose fates the model is built to compare: homeless and
unemployed (health zero, unhireable); renting and unemployed (one real-estate
day from homelessness); renting and employed (housed but assetless, one shock
from the trap); owning and employed (secure, accumulating); owning and
unemployed (asset-rich, perhaps cash-poor, exposed to the illiquidity death);
the \textbf{rentier} who lives off rents and liquid assets without working; the
\textbf{entrepreneur} who has risked a home to run a firm; and the
\textbf{indigenous} mine owner at the top. The through-line is \textbf{why you
need savings}: with annual hiring, annual rent, and monthly food, only a cash
buffer carries an agent across a gap in income --- and running it down is the
common path into the trap.

\subsection{The initial seed}
\label{sec:seed}

The economy begins stable, fully employed, and at level~0 --- automation has not
started and must be ignited from within. Every individual starts housed (owning
or renting), with a job, with health and happiness at ten, and with the starting
education its job requires. Wealth is drawn at random across three strata ---
a \emph{bourgeoisie} owning their home plus extra real estate and debt holdings,
a \emph{middle} owning only their home, and a \emph{poor} owning nothing and so
renting --- but this is asset ownership, not shelter: nobody is homeless at the
start. The \textbf{indigenous} own all mining land and the mines, which makes
them the economy's primary source of capital from turn one; they capitalise the
bank, collect its dividend, and lend to entrepreneurs, so financial power is
triply concentrated in them at the outset. Critically, there are regular
engineers but \textbf{no advanced engineers}: that class starts empty, which is
what forces ignition.

\subsection{The automation shock, stepped through}

With only level-0 machines and no advanced engineers, the only thing a business
can do with profit at first is buy more level-0 machines --- so the economy can
only \emph{expand}, hiring more operators for more output. It cannot yet
automate. Ignition is endogenous and runs in sequence: an agent acquires enough
education to become the first \textbf{advanced engineer}; a tech company hires it
and builds a level-1 machine; businesses buy the new generation, which makes the
same or more output with fewer operators; and the fire-or-redeploy choice
finally has something to bite on. The shock's \emph{timing} is therefore set by
how fast the education system produces that first advanced engineer --- so
university funding and mix are the master lever on \emph{when} the shock fires,
not merely on how unequal its aftermath is.

From there the model should, if the micro-rules suffice, generate the whole
arc without being told to: rising supply and falling employment together;
unsold goods and falling prices; a shrinking money supply as lending dries up;
a rising real burden of fixed debts; distressed firms ratcheting up offered
rates and holders dumping marked-down debt; house and mine-land fire-sales as
banks and folding firms raise cash; a bank base contracting through its own
fire-sale; and a government whose rescue succeeds or fails at the verdict of the
debt market. And beneath it all, the dark stabiliser of
\S\ref{sec:birthdeath}: a death-wave that cancels debt and hands houses to the
bottom, letting the system find a floor through attrition.

\hyp A single declining parameter --- operators per machine --- is sufficient,
under boundedly-rational micro-rules and stock-flow-consistent money, to produce
both of Fisher's movements: a deflation in goods and a debt-deflation in assets,
with structural unemployment between them. The experiment is to run the economy
with this change permitted and again with it forbidden, and compare posted
prices, new lending, and deaths.

\subsection{Resolved questions, and what remains to calibrate}
\label{sec:open}

The design is closed; every structural question has a decision. Resolved earlier:
mining is labour-free with no miner career; the machine market trades on job day;
mine-land yield is monthly, unbounded up to a mine's finite stock; rent is annual
and paid upfront, so housing is secure within the year (no mid-year
homelessness); fertility gates are fixed and the voluntary choice follows the
agent's objective; the lookahead horizon follows an agent's time-preference; the
bank's lending base is marked to the independent mine-land market
(\S\ref{sec:bank}); agents choose their education from last year's job offerings;
a margin call is met by selling liquid assets before any foreclosure; and the
glut branch is just the protocol's ``cut the price of what went unsold,'' so the
symmetric price-fall is already built in. Resolved here:

\begin{itemize}
  \item \textbf{Secured-first in a childless estate}: the bank seizes its lien'd
  collateral first; only unsecured owed debts (and any deficiency) are written
  off, so the homeless receive houses net of any lien and the bank bears a loss
  only when collateral is worth less than the loan. A sub-ten-unit remainder
  escheats to the government with the liquid assets; an heir inherits the net
  estate only, never liabilities.
  \item \textbf{No forced deleveraging}: the ten-times ceiling governs only
  \emph{new} lending (it stops as the base shrinks); existing loans are
  disciplined solely by per-loan margin calls against their own collateral ---
  the bank never recalls the book en masse.
  \item \textbf{Labour-force inflow and retraining}: births are the only inflow
  (no immigration); a career switch resets education to zero and re-education then
  proceeds at the ordinary rate, so there is no separate retraining speed.
  \item \textbf{Scripted-agent objectives are deliberately the implementer's
  choice} --- what the government optimises and an entrepreneur's
  reinvest-versus-extract split are left open by design, not by omission.
\end{itemize}

\noindent What remains is \emph{numeric calibration}, not design: the required
cash gift $G$ (a newborn's endowment, default about a year of subsistence food),
the bureaucrat support ratio $\beta$ (default one at generation~0), the
raw-material recipe coefficients per generation, and the usual constants --- the
loan-to-value haircut, the ten-times leverage cap, the sales-tax rate, the
fire-sale discount, the total raw-material stock, and the $t_0$ counts. Each is a
knob to set and sweep, with no design decision outstanding.

\end{document}
